\documentclass[10pt,a4paper]{amsart}
\usepackage[margin=19mm]{geometry}
\usepackage{amsmath,amssymb,mathtools}
\usepackage[scr=rsfs,cal=euler]{mathalfa}
\usepackage{xfrac}
\usepackage[unicode,hidelinks,bookmarks=false]{hyperref}
\newcommand{\ie}{i.e.\ }
\newcommand{\cf}{cf.\ }
\newcommand{\resp}{respectively\ }
\newcommand{\Subsection}{\subsection}
\providecommand{\nobreakdash}{\nobreak\hbox{-}}
\setlength{\emergencystretch}{2em}
\multlinegap0pt
\usepackage{amssymb}
\usepackage{xcolor}
\newtheorem{lm}{Lemma}
\newtheorem{thm}{Theorem}
 \newcommand{\FF}{\mathbb F}
\title{On additive MDS codes with linear projections}
\author{Sam Adriaensen, Simeon Ball}
\date{Source: arXiv:2209.09767v1 (CC BY 4.0)}
\begin{document}
\maketitle

\noindent\emph{Source and licence.} On additive MDS codes with linear projections. Version arXiv:2209.09767v1 (CC BY 4.0), distributed by arXiv under the Creative Commons Attribution 4.0 International licence, http://creativecommons.org/licenses/by/4.0/, as recorded in the arXiv licence metadata for this exact version and shown on the abstract page https://arxiv.org/abs/2209.09767v1. Full licence notice: \texttt{licenses/arxiv-2209.09767-CC-BY-4.0.txt}.\par\smallskip

\noindent\textit{Editorial scope.} Counted unit: the article's thm thm (source0.tex lines 969--972). The article's complete original proof of it is delivered with it (source0.tex lines 974--1020, one \texttt{proof} environment of 47 lines). No mathematics is written, repaired or re-derived: the statement and the proof are the article's own lines, in the article's own notation, and no retained line is substituted or re-flowed.  Retained uncounted as the source-local context the proof needs: lines 919--923.  Nothing was omitted from any retained span, so the package carries no omission record.  No retained line contains a citation, so the wrapper carries no bibliography and no bibliography entry.  No retained line contains a figure environment, an image-inclusion command or drawing code, so no image is delivered and the package declares no asset dependency.  Editorial formatting only, in addition: an AMS \texttt{amsart} preamble that restates the article's own declarations and macros used by the retained lines, namely the environments \texttt{lm}, \texttt{thm} on separate counters, together with the standard packages those lines need.  The retained environments print in this document as Lm 1, Theorem 1 in order of appearance, whereas the article prints the same items with the numbering of its own full document.  The complete native source of the article as distributed in the arXiv e-print for this exact version is bundled unchanged as \texttt{sources/130972/source0.tex}.
\begin{lm}
 \label{LmProp}
 Suppose that for a given value of $m$, the only pairs of $\FF_q$-linearised polynomials of $\FF_{q^h}$ satisfying property $(Prop_m)$ are monomials.
 Then every $\FF_q$-linear $(n,q^{3h},n-2)_{q^h}$ code $C$, with $n \geq 3 + m$ which has at least three projections that are equivalent to linear codes, is itself equivalent to a linear code.
\end{lm}

\begin{thm}
 Let $C$ be an $\FF_q$-linear $(n,q^6,n-2)_{q^2}$ MDS code over $\FF_{q^2}$, with $n \geq q + 3$.
 If $C$ has at least three coordinates from which the projections are equivalent to linear codes, then $C$ itself is equivalent to a linear code.
\end{thm}

\begin{proof}
By Lemma \ref{LmProp}, it suffices to show that if $f$ and $g$ satisfy $(Prop_q)$, then $f$ and $g$ are monomials.
Thus, suppose that
\[
 a_j f(b_j f^{-1} (X)) \equiv g(c_j g^{-1}(X))
\]
for $j \in [1,q]$, with all $a_j$, all $b_j$, and all $c_j$ non-zero and distinct.
The Dickson matrix of $f$ equals
\(
 M_f = \begin{pmatrix}
  f_0 & f_1 \\ f_1^q & f_0^q
 \end{pmatrix}
\).
Thus,
\[
 M_{f^{-1}} = M_f^{-1} = \frac 1 {f_0^{q+1}-f_1^{q+1}} \begin{pmatrix}
 f_0^q & -f_1 \\ -f_1^q & f_0
 \end{pmatrix}.
\]
Therefore,
\[
 f^{-1}(X) = \frac{f_0^q X - f_1 X^q}{f_0^{q+1}-f_1^{q+1}}.
\]
Using the fact that 
\(
 f_0^{q+1} - f_1^{q+1} \in \FF_q^*
\), this means that
\begin{align}
 a_j f(b_j f^{-1}(X))
 \equiv \frac{a_j}{f_0^{q+1} - f_1^{q+1}} \left( (f_0^{q+1} b_j - f_1^{q+1} b_j^q) X + f_1 f_0 (b_j^q-b_j) X^q \right).
\end{align}
We can do the same for $g(c_j g^{-1} (X))$.
Looking at the coefficient of $X^q$ tells us that
\[
 \frac{a_j}{f_0^{q+1} - f_1^{q+1}} f_1 f_0 (b_j^q - b_j)
 = \frac 1 {g_0^{q+1} - g_1^{q+1}} g_1 g_0 (c_j^q - c_j).
\]
If we raise both sides of this equation to the power $q-1$, this tells us
\[
 (f_0 f_1)^{q-1} a_j^{q-1} = (g_1 g_0)^{q-1}.
\]
This holds for $q$ distinct non-zero values of $a_j$, so this is only possible if
\(
 (f_0 f_1)^{q-1} = (g_0 g_1)^{q-1} = 0,
\)
which implies that $f$ and $g$ are monomials.
\end{proof}

\end{document}
